% Autour du nombre e : factorielles, moyennes, irrationalité
% Compile avec pdflatex, xelatex ou lualatex.
\documentclass[11pt,a4paper]{article}

\usepackage{iftex}
\ifPDFTeX
  \usepackage[utf8]{inputenc}
  \usepackage[T1]{fontenc}
  \IfFileExists{lmodern.sty}{\usepackage{lmodern}}{}
  \usepackage[french]{babel}
\else
  \usepackage{fontspec}
  \usepackage{polyglossia}
  \setdefaultlanguage{french}
\fi
\usepackage[margin=2.6cm]{geometry}
\usepackage{amsmath,amssymb,amsthm,mathtools}
\usepackage{enumitem}
\usepackage{microtype}
\usepackage[hidelinks]{hyperref}

% ---------- Environnements ----------
\theoremstyle{plain}
\newtheorem{theoreme}{Théorème}[section]
\newtheorem{proposition}[theoreme]{Proposition}
\newtheorem{lemme}[theoreme]{Lemme}
\newtheorem{corollaire}[theoreme]{Corollaire}
\theoremstyle{definition}
\newtheorem{definition}[theoreme]{Définition}
\newtheorem{exemple}[theoreme]{Exemple}
\newtheorem{exercice}{Exercice}
\theoremstyle{remark}
\newtheorem{remarque}[theoreme]{Remarque}

% ---------- Notations ----------
\newcommand{\N}{\mathbb{N}}
\newcommand{\Z}{\mathbb{Z}}
\newcommand{\Q}{\mathbb{Q}}
\newcommand{\R}{\mathbb{R}}
\newcommand{\dd}{\,\mathrm{d}}

\title{Autour du nombre $e$\\[2pt]
\large Factorielles, inégalités de moyennes, irrationalité}
\author{Michel Hua}
\date{Octobre 2026}

\begin{document}
\maketitle

\begin{abstract}
On reprend une feuille de problèmes de terminale consacrée à l'exponentielle, dont les trois parties tournent autour du nombre $e$. L'encadrement $2\leq(1+\frac1n)^n\leq e$ donne un premier encadrement de $n!$, que l'on pousse jusqu'à la formule de Stirling avec un reste explicite. L'inégalité $\ln x\leq x-1$ donne l'inégalité arithmético-géométrique, que l'on replace dans la famille des moyennes d'ordre $p$. Enfin, les sommes partielles de $\sum 1/k!$ fournissent deux suites adjacentes de limite $e$, la formule de Taylor avec reste intégral en donne la vitesse, et l'on en déduit l'irrationalité de $e$ et un développement asymptotique de $n!\,e$ modulo $1$.
\end{abstract}

\paragraph{Leçons concernées.} 224 (Exemples de développements asymptotiques de suites et de fonctions), 223 (Suites numériques), 229 (Fonctions monotones, fonctions convexes), 218 (Formules de Taylor), 230 (Séries numériques).

\paragraph{Développements possibles.} Formule de Stirling avec reste (théorème \ref{thm:stirling}) ; monotonie des moyennes d'ordre $p$ (théorème \ref{thm:moyennes}) ; reste de la série de $e$ et irrationalité de $e$ (proposition \ref{prop:reste} et théorème \ref{thm:irr}).

\paragraph{Notations.} On note $u_n\sim v_n$ l'équivalence de suites et $O(\cdot)$ la domination usuelle. Les fonctions $\ln$ et $\exp$ sont supposées connues, avec $e=\exp(1)$.

\tableofcontents

% ==========================================================
\section{Factorielles et formule de Stirling}
% ==========================================================

\subsection{Un premier encadrement}

\begin{lemme}\label{lem:2e}
Pour tout $n\geq1$, on a $2\leq\big(1+\frac1n\big)^n<e$.
\end{lemme}

\begin{proof}
Par la formule du binôme, $\big(1+\frac1n\big)^n\geq1+n\cdot\frac1n=2$ (inégalité de Bernoulli). Par concavité de $\ln$, on a $\ln(1+x)\leq x$ pour $x>-1$, avec égalité seulement en $x=0$ ; donc $n\ln\big(1+\frac1n\big)<1$, et l'on applique $\exp$, qui est strictement croissante.
\end{proof}

\begin{proposition}\label{prop:encadrement}
Soit $u_n=\dfrac{n^n}{n!}$ pour $n\geq1$. Alors $\dfrac{u_{n+1}}{u_n}=\big(1+\frac1n\big)^n$, la suite $(u_n)$ est strictement croissante et, pour tout $n\geq1$,
\[
  2^{n-1}\leq u_n\leq e^{n-1},
  \qquad\text{c'est-à-dire}\qquad
  e\Big(\frac ne\Big)^n\leq n!\leq 2\Big(\frac n2\Big)^n .
\]
De plus $n!\leq\big(\frac n2\big)^n$ pour tout $n\geq6$.
\end{proposition}

\begin{proof}
On a $\dfrac{u_{n+1}}{u_n}=\dfrac{(n+1)^{n+1}}{(n+1)!}\cdot\dfrac{n!}{n^n}=\dfrac{(n+1)^n}{n^n}$. Par le lemme \ref{lem:2e}, ce quotient est dans $[2,e[$, donc $(u_n)$ est strictement croissante. Comme $u_1=1$, le produit télescopique $u_n=\prod_{k=1}^{n-1}u_{k+1}/u_k$ donne $2^{n-1}\leq u_n\leq e^{n-1}$, ce qui se réécrit comme annoncé.

Pour la dernière assertion, on a $u_6=\frac{46656}{720}=64{,}8\geq2^6$, et $u_{n+1}\geq2u_n$ donne $u_n\geq2^n$ pour $n\geq6$ par récurrence. (L'inégalité est fausse pour $n=5$ : $u_5=\frac{3125}{120}<32$.)
\end{proof}

\begin{remarque}
La minoration $n!\geq(n/e)^n$, un peu plus faible que celle de la proposition, s'obtient en une ligne : $e^n=\sum_{k\geq0}\frac{n^k}{k!}\geq\frac{n^n}{n!}$.
\end{remarque}

\begin{corollaire}\label{cor:croissances}
Pour tout réel $a>0$, $\dfrac{n!}{a^n}\to+\infty$ ; autrement dit $a^n=o(n!)$.
\end{corollaire}

\begin{proof}
Par la proposition \ref{prop:encadrement}, $\dfrac{n!}{a^n}\geq e\Big(\dfrac{n}{ea}\Big)^n$, et $\dfrac{n}{ea}\geq2$ dès que $n\geq2ea$, donc le minorant tend vers $+\infty$.
\end{proof}

\subsection{La formule de Stirling}

L'encadrement précédent a le bon ordre de grandeur exponentiel, $\ln n!=n\ln n+O(n)$, mais laisse un facteur $(e/2)^n$ d'incertitude. On le raffine en étudiant la suite
\[
  w_n=\ln n!-\Big(n+\frac12\Big)\ln n+n,\qquad n\geq1 .
\]

\begin{lemme}\label{lem:wn}
Pour tout $n\geq1$,
\[
  \frac1{3(2n+1)^2}\leq w_n-w_{n+1}\leq\frac1{12n(n+1)}=\frac1{12}\Big(\frac1n-\frac1{n+1}\Big).
\]
En particulier $w_n-w_{n+1}=\dfrac1{12n^2}+O\Big(\dfrac1{n^3}\Big)$.
\end{lemme}

\begin{proof}
On calcule
\[
  w_n-w_{n+1}=-\ln(n+1)-\Big(n+\frac12\Big)\ln n+\Big(n+\frac32\Big)\ln(n+1)-1
  =\Big(n+\frac12\Big)\ln\Big(1+\frac1n\Big)-1 .
\]
Posons $t=\frac1{2n+1}\in\,]0,1[$. Alors $\frac{1+t}{1-t}=\frac{2n+2}{2n}=1+\frac1n$ et $n+\frac12=\frac1{2t}$. Le développement en série entière $\ln\frac{1+t}{1-t}=2\sum_{k\geq0}\frac{t^{2k+1}}{2k+1}$, valable pour $|t|<1$, donne
\[
  w_n-w_{n+1}=\frac1{2t}\ln\frac{1+t}{1-t}-1=\sum_{k\geq1}\frac{t^{2k}}{2k+1}.
\]
Cette somme est minorée par son premier terme $\frac{t^2}3$, et majorée par $\frac13\sum_{k\geq1}t^{2k}=\frac{t^2}{3(1-t^2)}$. Or $\frac{t^2}{1-t^2}=\frac1{(2n+1)^2-1}=\frac1{4n(n+1)}$, d'où l'encadrement. Les deux bornes valent $\frac1{12n^2}+O(\frac1{n^3})$.
\end{proof}

Pour identifier la limite, on utilise les intégrales de Wallis, étudiées en détail dans la fiche \texttt{lecons/236-calcul-integrales/calcul-integral.tex}. On rappelle l'énoncé et une preuve courte.

\begin{lemme}[Wallis]\label{lem:wallis}
Soit $W_n=\int_0^{\pi/2}\sin^nt\dd t$. Alors $W_n\sim\sqrt{\dfrac{\pi}{2n}}$, et
\[
  W_{2p}=\frac{(2p)!}{4^p\,(p!)^2}\cdot\frac\pi2,
  \qquad\text{donc}\qquad
  \frac{(2p)!}{4^p\,(p!)^2}\sim\frac1{\sqrt{\pi p}} .
\]
\end{lemme}

\begin{proof}
Une intégration par parties ($\sin^{n+2}=\sin^{n+1}\cdot\sin$) donne $(n+2)W_{n+2}=(n+1)W_n$. Avec $W_0=\frac\pi2$, on obtient la formule pour $W_{2p}$, et la suite $(n+1)W_{n+1}W_n$ est constante, égale à $W_1W_0=\frac\pi2$. Comme $0\leq\sin t\leq1$ sur $[0,\frac\pi2]$, la suite $(W_n)$ est décroissante et strictement positive, donc
\[
  \frac{n+1}{n+2}=\frac{W_{n+2}}{W_n}\leq\frac{W_{n+1}}{W_n}\leq1,
\]
d'où $W_{n+1}\sim W_n$, puis $W_n^2\sim\frac{\pi}{2(n+1)}\sim\frac\pi{2n}$. Enfin $\frac{(2p)!}{4^p(p!)^2}=\frac2\pi W_{2p}\sim\frac2\pi\sqrt{\frac{\pi}{4p}}=\frac1{\sqrt{\pi p}}$.
\end{proof}

\begin{theoreme}[Stirling]\label{thm:stirling}
On a $n!\sim\sqrt{2\pi n}\,\Big(\dfrac ne\Big)^n$. Plus précisément, pour tout $n\geq1$, il existe $\theta_n\in\,]0,1]$ tel que
\[
  n!=\sqrt{2\pi n}\,\Big(\frac ne\Big)^n e^{\theta_n/(12n)} .
\]
\end{theoreme}

\begin{proof}
Par le lemme \ref{lem:wn}, $(w_n)$ est strictement décroissante et $\big(w_n-\frac1{12n}\big)$ est croissante. Comme $w_n-\frac1{12n}\leq w_n$, ces deux suites sont adjacentes et convergent vers une même limite $\ell$. En sommant l'encadrement du lemme \ref{lem:wn} de $k=n$ à $+\infty$, on obtient $0<w_n-\ell\leq\frac1{12n}$, c'est-à-dire $n!=e^{\ell}\sqrt n\,(n/e)^n e^{\theta_n/(12n)}$ avec $\theta_n\in\,]0,1]$.

Il reste à montrer que $C=e^\ell$ vaut $\sqrt{2\pi}$. Puisque $n!\sim C\sqrt n\,(n/e)^n$,
\[
  \frac{(2p)!}{4^p(p!)^2}\sim\frac{C\sqrt{2p}\,(2p/e)^{2p}}{4^p\,C^2p\,(p/e)^{2p}}=\frac{\sqrt2}{C\sqrt p},
\]
et la comparaison avec le lemme \ref{lem:wallis} donne $\frac{\sqrt2}C=\frac1{\sqrt\pi}$, soit $C=\sqrt{2\pi}$.
\end{proof}

\begin{remarque}
\begin{enumerate}[nosep]
\item La décroissance de $(w_n)$ donne aussi l'encadrement $\sqrt{2\pi n}\,(n/e)^n\leq n!\leq e\sqrt n\,(n/e)^n$, puisque $\ell\leq w_n\leq w_1=1$. Il est à comparer avec la proposition \ref{prop:encadrement}.
\item Le lemme \ref{lem:wn} montre que $w_n-\ell=\sum_{k\geq n}(w_k-w_{k+1})=\frac1{12n}+O(\frac1{n^2})$. On obtient ainsi le début du développement asymptotique $\ln n!=n\ln n-n+\frac12\ln(2\pi n)+\frac1{12n}+O(\frac1{n^2})$. En poussant le calcul, on obtient la série de Stirling, dont les coefficients s'expriment à l'aide des nombres de Bernoulli et qui diverge pour tout $n$.
\item Pour $n=10$ : $10!=3\,628\,800$, et $\sqrt{20\pi}\,(10/e)^{10}\approx3\,598\,696$, soit une erreur relative d'environ $0{,}83\,\%\approx\frac1{120}$.
\end{enumerate}
\end{remarque}

% ==========================================================
\section{Inégalités entre moyennes}
% ==========================================================

Soit $n\geq2$ et $a=(a_1,\dots,a_n)$ des réels \emph{strictement positifs}. On note
\[
  m=\frac1n\sum_{i=1}^na_i,\qquad
  g=\Big(\prod_{i=1}^na_i\Big)^{1/n},\qquad
  \frac1h=\frac1n\sum_{i=1}^n\frac1{a_i}
\]
leurs moyennes arithmétique, géométrique et harmonique.

\begin{theoreme}\label{thm:hgm}
On a $h\leq g\leq m$. Chacune de ces deux inégalités est une égalité si et seulement si $a_1=\cdots=a_n$.
\end{theoreme}

\begin{proof}[Première preuve (méthode de la feuille)]
La fonction $\varphi(x)=x-1-\ln x$ vérifie $\varphi'(x)=1-\frac1x$, donc $\varphi$ décroît sur $]0,1]$ et croît sur $[1,+\infty[$. Elle a un minimum strict en $1$, où elle vaut $0$. Ainsi $\ln x\leq x-1$ pour tout $x>0$, avec égalité si et seulement si $x=1$. En appliquant cette inégalité à $x=a_i/m$ et en sommant,
\[
  n\ln\frac gm=\sum_{i=1}^n\ln\frac{a_i}m\leq\sum_{i=1}^n\Big(\frac{a_i}m-1\Big)=0,
\]
d'où $g\leq m$, avec égalité si et seulement si $a_i=m$ pour tout $i$. On applique ensuite ce résultat aux réels $1/a_i$, dont la moyenne géométrique vaut $1/g$ et la moyenne arithmétique $1/h$. On obtient $\frac1g\leq\frac1h$, soit $h\leq g$, avec le même cas d'égalité.
\end{proof}

\begin{proof}[Seconde preuve (convexité)]
La fonction $\ln$ est strictement concave sur $]0,+\infty[$, car $\ln''(x)=-1/x^2<0$. L'inégalité de Jensen donne $\frac1n\sum\ln a_i\leq\ln\big(\frac1n\sum a_i\big)$, c'est-à-dire $\ln g\leq\ln m$, et la stricte concavité donne le cas d'égalité. On conclut pour $h$ comme ci-dessus.
\end{proof}

\begin{remarque}
La première preuve est en fait celle de Jensen pour $\ln$, la droite $y=x-1$ étant la tangente au graphe de $\ln$ en $1$. Remplacer $1$ par $m$ revient à se placer au point où l'on applique Jensen. C'est la preuve « par la tangente » de l'inégalité de Jensen, qui vaut pour toute fonction concave dérivable.
\end{remarque}

\begin{exemple}[Exercice 2 de la feuille]\label{ex:4-10-25}
Pour $(a_1,a_2,a_3)=(4,10,25)$ : $m=\frac{39}3=13$, $g=\sqrt[3]{1000}=10$ et $\frac1h=\frac13\big(\frac14+\frac1{10}+\frac1{25}\big)=\frac13\cdot\frac{39}{100}=\frac{13}{100}$, donc $h=\frac{100}{13}\approx7{,}69$. On a bien $h<g<m$, et de plus $m\,h=100=g^2$.

Cette dernière égalité n'est pas un hasard : $4,10,25$ est une progression géométrique de raison $\frac52$. Pour $(a,ar,ar^2)$ avec $a,r>0$, on a $g=ar$, $m=\frac{a(1+r+r^2)}3$ et $\frac1h=\frac{1+r+r^2}{3ar^2}$, donc $mh=a^2r^2=g^2$. Pour $n=2$, l'identité $mh=g^2$ vaut pour tout couple, car $\frac{a+b}2\cdot\frac{2ab}{a+b}=ab$. En revanche, pour $(1,1,4)$, on a $m=2$, $g=\sqrt[3]4$ et $h=\frac43$, donc $mh=\frac83\neq g^2=\sqrt[3]{16}\approx2{,}52$.
\end{exemple}

\begin{exemple}\label{ex:croissance-e}
La suite $e_n=\big(1+\frac1n\big)^n$ est strictement croissante. En effet, appliquons $g<m$ aux $n+1$ réels $1+\frac1n,\dots,1+\frac1n,1$, qui ne sont pas tous égaux. On obtient
\[
  e_n^{1/(n+1)}<\frac{n(1+\frac1n)+1}{n+1}=1+\frac1{n+1},
\]
c'est-à-dire $e_n<e_{n+1}$. Avec le lemme \ref{lem:2e}, on retrouve que $(e_n)$ converge, en croissant, vers une limite $\leq e$ ; cette limite vaut $e$ puisque $n\ln(1+\frac1n)\to1$.
\end{exemple}

\subsection{Moyennes d'ordre \texorpdfstring{$p$}{p}}

\begin{definition}
Pour $p\in\R^*$, on pose $M_p(a)=\Big(\dfrac1n\sum_{i=1}^na_i^p\Big)^{1/p}$, et $M_0(a)=g$.
\end{definition}

Ainsi $M_1=m$, $M_{-1}=h$, et $M_2$ est la moyenne quadratique. On a aussi $M_{-p}(a)=1/M_p(1/a)$, où $1/a=(1/a_1,\dots,1/a_n)$.

\begin{theoreme}\label{thm:moyennes}
L'application $p\mapsto M_p(a)$ est continue et croissante sur $\R$. Elle est strictement croissante si les $a_i$ ne sont pas tous égaux, et constante sinon. De plus $M_p(a)\to\max a_i$ quand $p\to+\infty$ et $M_p(a)\to\min a_i$ quand $p\to-\infty$.
\end{theoreme}

\begin{proof}
Si les $a_i$ sont tous égaux à $c$, alors $M_p(a)=c$ pour tout $p$. Supposons-les non tous égaux.

\emph{Cas $0<p<q$.} La fonction $x\mapsto x^{q/p}$ est strictement convexe sur $]0,+\infty[$ puisque $q/p>1$. Par Jensen appliqué aux réels $b_i=a_i^p$, qui ne sont pas tous égaux,
\[
  M_p(a)^q=\Big(\frac1n\sum b_i\Big)^{q/p}<\frac1n\sum b_i^{q/p}=M_q(a)^q ,
\]
et l'on prend la racine $q$-ième.

\emph{Cas $p<q<0$.} On a $0<-q<-p$, donc, par le cas précédent appliqué à $1/a$, $M_{-q}(1/a)<M_{-p}(1/a)$, et en passant aux inverses $M_p(a)<M_q(a)$.

\emph{Cas où $p$ ou $q$ est nul.} Pour $q>0$, le théorème \ref{thm:hgm} appliqué aux $a_i^q$ donne $g^q<M_q(a)^q$, soit $M_0<M_q$. Pour $p<0$, il donne de même $g^p<M_p(a)^p$, et comme $x\mapsto x^{1/p}$ est décroissante, $M_p(a)<g=M_0(a)$. Le cas $p<0<q$ s'en déduit par transitivité.

\emph{Continuité en $0$.} Pour $p\to0$, on a $a_i^p=e^{p\ln a_i}=1+p\ln a_i+O(p^2)$, donc
\[
  \ln M_p(a)=\frac1p\ln\Big(1+\frac pn\sum_i\ln a_i+O(p^2)\Big)=\frac1n\sum_i\ln a_i+O(p)\longrightarrow\ln g .
\]
Ailleurs, la continuité est claire.

\emph{Limites.} Soit $A=\max a_i$. Pour $p>0$, on a $n^{-1/p}A\leq M_p(a)\leq A$, et $n^{-1/p}\to1$. On en déduit le cas $p\to-\infty$ par $M_{-p}(a)=1/M_p(1/a)$.
\end{proof}

\begin{remarque}
En particulier $\min a_i\leq h\leq g\leq m\leq M_2\leq\max a_i$. Sur l'exemple \ref{ex:4-10-25}, $M_2=\sqrt{247}\approx15{,}7$.
\end{remarque}

% ==========================================================
\section{Deux suites adjacentes de limite \texorpdfstring{$e$}{e}}
% ==========================================================

Pour $n\geq1$, on pose
\[
  u_n=\sum_{k=0}^n\frac1{k!},\qquad v_n=u_n+\frac1{n\cdot n!} .
\]
(Dans cette section, $u_n$ ne désigne plus la suite de la section 1.)

\begin{proposition}\label{prop:adjacentes}
Les suites $(u_n)$ et $(v_n)$ sont adjacentes : $(u_n)$ est strictement croissante, $(v_n)$ est strictement décroissante et $v_n-u_n\to0$. Leur limite commune $\ell$ vérifie donc $u_n<\ell<v_n$ pour tout $n\geq1$.
\end{proposition}

\begin{proof}
On a $u_{n+1}-u_n=\frac1{(n+1)!}>0$ et $v_n-u_n=\frac1{n\cdot n!}\to0$. Enfin
\[
  v_{n+1}-v_n=\frac1{(n+1)!}+\frac1{(n+1)(n+1)!}-\frac1{n\cdot n!}
  =\frac{n(n+1)+n-(n+1)^2}{n(n+1)\,(n+1)!}=-\frac1{n(n+1)\,(n+1)!}<0 .
\]
Les inégalités strictes $u_n<\ell<v_n$ viennent de la stricte monotonie.
\end{proof}

\begin{exemple}
On a $u_1=2$, $v_1=3$, $u_2=\frac52$, $v_2=\frac{11}4$, $u_3=\frac83$ et $v_3=\frac{49}{18}$. L'amplitude $v_n-u_n=\frac1{n\cdot n!}$ vaut $\frac1{600}$ pour $n=5$ et $\frac1{4320}<10^{-3}$ pour $n=6$. Avec $u_6=\frac{1957}{720}\approx2{,}71806$ et $v_6=\frac{11743}{4320}\approx2{,}71829$, on obtient $2{,}718<\ell<2{,}7183$. On conjecture $\ell=e$.
\end{exemple}

\begin{proposition}\label{prop:reste}
Pour tout $n\geq1$,
\[
  e-u_n=\frac1{n!}\int_0^1(1-t)^ne^t\dd t ,
  \qquad\text{donc}\qquad
  \frac1{(n+1)!}<e-u_n<\frac{e}{(n+1)!} .
\]
En particulier $u_n\to e$, donc $\ell=e$, et $e=\sum_{k\geq0}\frac1{k!}$.
\end{proposition}

\begin{proof}
C'est la formule de Taylor avec reste intégral appliquée à $\exp$ entre $0$ et $1$, à l'ordre $n$. Pour rester au niveau de la feuille, on peut la retrouver ainsi. Soit
\[
  f(x)=\Big(\sum_{k=0}^n\frac{x^k}{k!}\Big)e^{-x},\qquad x\in[0,1].
\]
Elle est dérivable, et
\[
  f'(x)=\Big(\sum_{k=1}^n\frac{x^{k-1}}{(k-1)!}-\sum_{k=0}^n\frac{x^k}{k!}\Big)e^{-x}=-\frac{x^n}{n!}e^{-x},
\]
car la somme télescope. Comme $f(0)=1$ et $f(1)=u_ne^{-1}$,
\[
  u_ne^{-1}-1=-\frac1{n!}\int_0^1x^ne^{-x}\dd x,
  \qquad\text{soit}\qquad
  e-u_n=\frac1{n!}\int_0^1x^ne^{1-x}\dd x ,
\]
et le changement de variable $t=1-x$ donne la formule. Pour $t\in\,]0,1[$, on a $1<e^t<e$, et $\int_0^1(1-t)^n\dd t=\frac1{n+1}$, d'où l'encadrement (strict, car l'intégrande est continue et les inégalités sont strictes sur $]0,1[$). Le majorant tend vers $0$.
\end{proof}

\begin{remarque}
\begin{enumerate}[nosep]
\item La proposition \ref{prop:adjacentes} donne un majorant plus fin que la proposition \ref{prop:reste} : $e-u_n<\frac1{n\cdot n!}$. La comparaison des deux minorants et majorants montre que $e-u_n\sim\frac1{(n+1)!}$ ; on précise ce point à l'exercice \ref{exo:sin}.
\item La convergence est très rapide : pour obtenir $d$ décimales, il suffit que $e/(n+1)!\leq10^{-d}$, et par Stirling le plus petit tel $n$ vérifie $n\ln n\sim d\ln10$. Par exemple $n=17$ suffit pour $15$ décimales, car $e/18!\approx4{,}2\cdot10^{-16}$.
\item Pour $x=-1$, la même méthode donne la série alternée $\frac1e=\sum_{k\geq0}\frac{(-1)^k}{k!}=\frac12-\frac1{2\cdot3}+\frac1{2\cdot3\cdot4}-\cdots$. C'est le développement de Pierce de $1/e$, de chiffres $q_i=i+1$, qui ne s'arrête pas ; cela redonne l'irrationalité de $e$ (voir la fiche \texttt{lecons/142-pgcd-ppcm/divisibilite-pgcd-congruences.tex}).
\end{enumerate}
\end{remarque}

\begin{theoreme}\label{thm:irr}
Le nombre $e$ est irrationnel.
\end{theoreme}

\begin{proof}
Supposons $e=\frac pq$ avec $p,q$ entiers, $q\geq1$. Par la proposition \ref{prop:adjacentes} et $\ell=e$,
\[
  u_q<\frac pq<u_q+\frac1{q\cdot q!} .
\]
En multipliant par $q\cdot q!>0$, on obtient $0<p\cdot q!-q\cdot q!\,u_q<1$. Or $q!\,u_q=\sum_{k=0}^q\frac{q!}{k!}$ est un entier, donc $p\cdot q!-q\cdot q!\,u_q$ est un entier strictement compris entre $0$ et $1$. C'est absurde.
\end{proof}

\begin{remarque}
\begin{enumerate}[nosep]
\item Autrement dit, si $e$ était rationnel de dénominateur $q$, le réel $q!\,(e-u_q)$ serait un entier. Or il est dans $\big]0,\frac1q\big[$. Le principe est que $e$ est approché par des rationnels $u_q$, de dénominateur divisant $q!$, beaucoup mieux que ne le permettrait un rationnel fixé.
\item La même idée, appliquée simultanément aux séries de $e$ et de $e^{-1}$, montre que $e^2$ est irrationnel et même non quadratique (Liouville, 1840) : si $e^2=\frac pq$, on écrit $qe=pe^{-1}$. Hermite a démontré en 1873 que $e$ est transcendant. On ne démontre pas ces résultats ici.
\end{enumerate}
\end{remarque}

% ==========================================================
\section{Exercices}
% ==========================================================

\begin{exercice}\label{exo:bernoulli}
Montrer que $(1+x)^n\geq1+nx$ pour tout $x\geq-1$ et tout $n\in\N$. Que donne la proposition \ref{prop:encadrement} pour $\ln n!$ ?
\end{exercice}

\begin{exercice}
Calculer $\displaystyle\lim_{n\to+\infty}\frac{n!}{n^n}$, $\displaystyle\lim_{n\to+\infty}\frac{(n!)^{1/n}}{n}$ et donner un équivalent de $\dbinom{2n}{n}$.
\end{exercice}

\begin{exercice}
Montrer que $\big(1+\frac1n\big)^n=e-\dfrac{e}{2n}+\dfrac{11e}{24n^2}+O\Big(\dfrac1{n^3}\Big)$.
\end{exercice}

\begin{exercice}
Soit $n\geq1$ et $g(x)=f(x)+\frac{x}{n!}$, où $f$ est la fonction de la preuve de la proposition \ref{prop:reste}. Montrer que $g$ est croissante sur $[0,1]$ et en déduire $e-\frac e{n!}\leq u_n\leq e$.
\end{exercice}

\begin{exercice}
Soit $x,y>0$ distincts. Montrer que la moyenne logarithmique $\frac{x-y}{\ln x-\ln y}$ s'intercale entre les moyennes géométrique et arithmétique : $\sqrt{xy}<\frac{x-y}{\ln x-\ln y}<\frac{x+y}2$.
\end{exercice}

\begin{exercice}\label{exo:sin}
On pose $r_n=n!\,(e-u_n)$. Montrer que $r_n=\dfrac1n+O\Big(\dfrac1{n^3}\Big)$, puis que $\sin(2\pi\,n!\,e)\sim\dfrac{2\pi}n$. En déduire la nature de la série $\sum\sin(2\pi\,n!\,e)$.
\end{exercice}

% ==========================================================
\section{Corrigés}
% ==========================================================

\paragraph{Exercice 1.} Par récurrence sur $n$ : si $(1+x)^n\geq1+nx$, alors, comme $1+x\geq0$, $(1+x)^{n+1}\geq(1+nx)(1+x)=1+(n+1)x+nx^2\geq1+(n+1)x$. Les encadrements de la proposition \ref{prop:encadrement} donnent, en prenant le logarithme, $n\ln n-n+1\leq\ln n!\leq n\ln n-(n-1)\ln2$, donc $\ln n!=n\ln n+O(n)$ et $\ln n!\sim n\ln n$.

\paragraph{Exercice 2.} On a $\frac{n!}{n^n}=\frac1{u_n}\leq\frac1{2^{n-1}}\to0$. Par Stirling, $\frac{(n!)^{1/n}}n=\frac1e(2\pi n)^{1/(2n)}e^{\theta_n/(12n^2)}\to\frac1e$, puisque $\frac{\ln(2\pi n)}{2n}\to0$. Enfin $\binom{2n}n=\frac{(2n)!}{(n!)^2}\sim\frac{4^n}{\sqrt{\pi n}}$ par le lemme \ref{lem:wallis}, ou directement par Stirling.

\paragraph{Exercice 3.} Avec $x=\frac1n$, on a $n\ln(1+x)=\frac1x\big(x-\frac{x^2}2+\frac{x^3}3+O(x^4)\big)=1-\frac x2+\frac{x^2}3+O(x^3)$. Donc
\[
  \Big(1+\frac1n\Big)^n=e\cdot\exp\Big(-\frac x2+\frac{x^2}3+O(x^3)\Big)
  =e\Big(1-\frac x2+\frac{x^2}3+\frac{x^2}8+O(x^3)\Big),
\]
et $\frac13+\frac18=\frac{11}{24}$. En particulier $e-\big(1+\frac1n\big)^n\sim\frac e{2n}$ : cette suite converge vers $e$ beaucoup plus lentement que $(u_n)$.

\paragraph{Exercice 4.} On a $g'(x)=\frac1{n!}\big(1-x^ne^{-x}\big)\geq0$ sur $[0,1]$, puisque $0\leq x^n\leq1$ et $e^{-x}\leq1$. Donc $g(1)\geq g(0)$, c'est-à-dire $u_ne^{-1}+\frac1{n!}\geq1$, soit $u_n\geq e-\frac e{n!}$. De même $f'\leq0$ donne $f(1)\leq f(0)$, soit $u_n\leq e$. Par encadrement, $u_n\to e$. C'est la méthode de la feuille d'origine, un peu moins précise que la proposition \ref{prop:reste}.

\paragraph{Exercice 5.} L'expression est symétrique en $x$ et $y$ ; on peut supposer $x>y$ et poser $s=\frac12\ln\frac xy>0$. En divisant par $\sqrt{xy}$, l'inégalité s'écrit
\[
  1\leq\frac{e^{s}-e^{-s}}{2s}=\frac{\operatorname{sh}s}s\leq\frac{e^{s}+e^{-s}}2=\operatorname{ch}s .
\]
La première inégalité vient de $\operatorname{sh}s=\sum_{k\geq0}\frac{s^{2k+1}}{(2k+1)!}\geq s$. La seconde équivaut à $\operatorname{th}s\leq s$, qui découle de $\operatorname{th}'=1-\operatorname{th}^2\leq1$ et $\operatorname{th}0=0$. Les inégalités sont strictes pour $s>0$.

\paragraph{Exercice 6.} On a $r_n=\sum_{k>n}\frac{n!}{k!}=\frac1{n+1}+\frac1{(n+1)(n+2)}+\rho_n$ avec
\[
  0<\rho_n=\sum_{k\geq n+3}\frac{n!}{k!}\leq\sum_{j\geq3}\frac1{(n+1)^j}=\frac1{n(n+1)^2}=O\Big(\frac1{n^3}\Big).
\]
Or $\frac1{n+1}=\frac1n-\frac1{n(n+1)}$ et $\frac1{(n+1)(n+2)}-\frac1{n(n+1)}=-\frac2{n(n+1)(n+2)}=O(\frac1{n^3})$, donc $r_n=\frac1n+O(\frac1{n^3})$. Plus précisément, on peut montrer que $r_n=\frac1n-\frac1{n^3}+O(\frac1{n^4})$.

Ensuite $n!\,e=n!\,u_n+r_n$ avec $n!\,u_n\in\N$, donc $\sin(2\pi\,n!\,e)=\sin(2\pi r_n)=2\pi r_n+O(r_n^3)=\frac{2\pi}n+O(\frac1{n^3})$. La série $\sum\sin(2\pi\,n!\,e)$ est à termes positifs à partir d'un certain rang (car $0<r_n<\frac1n$, d'après la proposition \ref{prop:adjacentes}) et son terme général est équivalent à celui de la série harmonique : elle diverge. En revanche, $\sum\big(\sin(2\pi\,n!\,e)-\frac{2\pi}n\big)$ converge absolument.

\begin{thebibliography}{9}
\bibitem{gourdon} X. Gourdon, \emph{Les maths en tête. Analyse}, Ellipses.
\bibitem{hw} G. H. Hardy, E. M. Wright, \emph{An Introduction to the Theory of Numbers}, Oxford University Press.
\bibitem{hlp} G. H. Hardy, J. E. Littlewood, G. Pólya, \emph{Inequalities}, Cambridge University Press.
\bibitem{rombaldi} J.-É. Rombaldi, \emph{Éléments d'analyse réelle}, EDP Sciences.
\end{thebibliography}

\paragraph{Source.} Feuille de problèmes de terminale S (2013), d'après le manuel \emph{Terracher}.

\end{document}
